%% ====================================================================
%%  Plantilla de evaluación
%% ====================================================================
\section{Plantilla de evaluación}\label{sec:plantilla}

A continuación, se presentan las plantillas de evaluación NPO
construidas para el análisis de las tecnologías proxy-caché y de
almacenamiento de archivos de gran tamaño. Las calificaciones siguen una
escala de 1 a 3, donde 1 es ``Muy limitado'', 2 es ``Aceptable'' y 3 es
``Excelente / Líder''.

En esta plantilla se presentan 6 indicadores que tienen un valor
ponderado con respecto a la puntuación final de cada una de las
herramientas; los indicadores que se tienen son evaluados con respecto a
las métricas que se habían definido en el formulario para la
caracterización de la organización. Cada una de estas métricas se
definen en la sección~\ref{sec:criterios}, criterios de evaluación, en donde se detalla la
importancia de cada una de estas en el contexto del proyecto.

A continuación, se le da un sentido a cada uno de los indicadores con
respecto a las métricas.

\begin{itemize}
\item \textbf{Rendimiento técnico (25\%):} Repositorio local de archivos grandes, priorización de tráfico académico, tiempos de descarga deficientes, optimización de ancho de banda, caché de archivos grandes, volúmenes de transferencia altos y tiempos de respuesta promedio.
\item \textbf{Escalabilidad (20\%):} Repositorio local de archivos grandes, priorización de tráfico académico (QoS), uso no misional de red, gobernanza del tráfico, QoS y priorización, ACL listas blancas/negras, escalabilidad de almacenamiento, políticas formales y diferenciar tráfico misional/no misional.
\item \textbf{Soporte y Mantenimiento (20\%):} Automatización de certificados SSL, falta de certificados SSL firmados, gestión de certificados digitales y adaptación a cambios.
\item \textbf{Adopción y Madurez (10\%):} Automatización de certificados SSL, proxy transparente/inverso/estándar e integración con servicios existentes.
\item \textbf{Observabilidad (10\%):} Uso no misional de red, métricas de uso de red y proxy, toma de decisiones basada en datos, monitoreo en tiempo real y alta disponibilidad, consumo de ancho de banda por usuario, estadísticas de tráfico por categoría y reportes.
\item \textbf{Adecuación al Contexto GRID (10\%):} Ajuste operativo de la herramienta en el contexto de operaciones del GRID y su escalabilidad.
\item \textbf{Riesgo Tecnológico (5\%):} Falta de certificados SSL firmados y sección para los estados proyectados para el año 2026 en la Matriz DAR.
\end{itemize}

\subsection{Cómo se calcula la puntuación de cada herramienta}\label{subsec:formalizacion}

Para que cualquier persona pueda reproducir o verificar los resultados,
en esta subsección se explica paso a paso cómo se pasa de las
calificaciones individuales de cada métrica a la puntuación DAR final de
una herramienta. El procedimiento consta de dos pasos: primero se calcula
el valor de cada indicador y después se combinan todos los indicadores en
una sola puntuación.

\subsubsection{Paso 1: obtener el valor de cada indicador}

Cada indicador agrupa varias métricas y todas ellas pesan lo mismo dentro
del indicador. Por eso, el valor del indicador es simplemente el
\textbf{promedio} de las calificaciones que la herramienta obtuvo en sus
métricas; es decir, se suman esas calificaciones y el resultado se divide
entre el número de métricas:

\[
  \text{Valor del indicador} \;=\;
  \frac{\text{suma de las calificaciones de sus métricas}}
       {\text{cantidad de métricas del indicador}}
\]

Como cada métrica se califica de 1 a 3, el promedio siempre queda también
entre 1 y 3. Esto es importante porque permite comparar indicadores entre
sí aunque unos agrupen más métricas que otros.

Según la lista anterior, cada indicador se calcula así:

\[
  \text{Rendimiento técnico} =
  \frac{\text{suma de sus 7 métricas}}{7}
  \qquad
  \text{Escalabilidad} =
  \frac{\text{suma de sus 9 métricas}}{9}
\]
\[
  \text{Soporte y mantenimiento} =
  \frac{\text{suma de sus 4 métricas}}{4}
  \qquad
  \text{Adopción y madurez} =
  \frac{\text{suma de sus 3 métricas}}{3}
\]
\[
  \text{Observabilidad} =
  \frac{\text{suma de sus 7 métricas}}{7}
\]

Por ejemplo, si una herramienta obtiene calificaciones de 3, 3, 2 y 2 en
las cuatro métricas de soporte y mantenimiento, el valor de ese indicador
es $(3+3+2+2) \div 4 = 2{,}5$.

Los dos indicadores restantes, adecuación al contexto GRID y riesgo
tecnológico, no se obtienen con este promedio, sino a partir de una
valoración del ajuste operativo de la herramienta en el GRID y de los
estados proyectados para el año 2026.

\subsubsection{Paso 2: obtener la puntuación DAR}

A diferencia de las métricas, los indicadores \textbf{no} pesan todos lo
mismo: cada uno tiene asignado un porcentaje según su importancia para el
proyecto. Por eso la puntuación DAR no es un promedio simple, sino un
\textbf{promedio ponderado}: el valor de cada indicador se multiplica por
su porcentaje y luego se suman todos los resultados.

\begin{align*}
  \text{DAR} \;=\;\; & 0{,}25 \times \text{Rendimiento técnico}
                 \;+\; 0{,}20 \times \text{Escalabilidad} \\[0.2em]
       &+\; 0{,}20 \times \text{Soporte y mantenimiento}
        \;+\; 0{,}10 \times \text{Adopción y madurez} \\[0.2em]
       &+\; 0{,}10 \times \text{Observabilidad}
        \;+\; 0{,}10 \times \text{Adecuación al contexto GRID} \\[0.2em]
       &+\; 0{,}05 \times \text{Riesgo tecnológico}
\end{align*}

Los factores $0{,}25$, $0{,}20$, $0{,}10$ y $0{,}05$ son los porcentajes
de la lista anterior escritos en forma decimal ($25\,\%$, $20\,\%$,
$10\,\%$ y $5\,\%$), y entre todos suman el $100\,\%$ de la evaluación.
Así, un buen resultado en rendimiento técnico influye cinco veces más en
la puntuación final que un buen resultado en riesgo tecnológico.

Como todos los indicadores están entre 1 y 3 y los porcentajes suman el
$100\,\%$, la puntuación DAR también queda entre 1 y 3: un valor cercano
a 3 indica una herramienta líder frente a los criterios priorizados por
el GRID, mientras que un valor cercano a 1 señala una alternativa muy
limitada para las necesidades del proyecto. Esta puntuación es la que se
reporta en la última columna de las plantillas y la que permite ordenar
las herramientas en el análisis DAR.

La Tabla~\ref{tab:plantilla-proxy} presenta la plantilla de análisis DAR empleada para evaluar
las herramientas de proxy-caché, con los criterios y sus respectivos
pesos.

\begin{table}[H]
\centering\footnotesize
\begin{tabular}{|>{\raggedright\arraybackslash\bfseries}p{2.4cm}|c|c|c|c|c|c|c|c|c|}
\hline
\rowcolor{gridhead} \thead[l]{Criterios/\\Herram.} & \thead{Frec.\\SMS} & \thead{Rend.\\(25\%)} & \thead{Esc.\\(20\%)} & \thead{Sop.\\(20\%)} & \thead{Adop.\\(10\%)} & \thead{Obs.\\(10\%)} & \thead{Adec.\\(10\%)} & \thead{Riesgo\\(5\%)} & \thead{DAR} \\ \hline
HAProxy & 0 &  &  &  &  &  &  &  &  \\ \hline
\rowcolor{gridrow} NGINX & 1-2 &  &  &  &  &  &  &  &  \\ \hline
Envoy & 1-2 &  &  &  &  &  &  &  &  \\ \hline
\rowcolor{gridrow} Varnish & 0 &  &  &  &  &  &  &  &  \\ \hline
Traefik & 0 &  &  &  &  &  &  &  &  \\ \hline
\rowcolor{gridrow} Caddy & 0 &  &  &  &  &  &  &  &  \\ \hline
Apache Traffic Server & 0 &  &  &  &  &  &  &  &  \\ \hline
\rowcolor{gridrow} Kong & 1 &  &  &  &  &  &  &  &  \\ \hline
Squid & 0 &  &  &  &  &  &  &  &  \\ \hline
\rowcolor{gridrow} Istio & 0 &  &  &  &  &  &  &  &  \\ \hline
Linkerd & 0 &  &  &  &  &  &  &  &  \\ \hline
\rowcolor{gridrow} K8Sidecar & 1 &  &  &  &  &  &  &  &  \\ \hline
Proxy-Service & 1 &  &  &  &  &  &  &  &  \\ \hline
\rowcolor{gridrow} nuster & 0 &  &  &  &  &  &  &  &  \\ \hline
\end{tabular}
\caption{Plantilla análisis DAR proxy-caché}
\label{tab:plantilla-proxy}
\end{table}

De igual manera, en la Tabla~\ref{tab:plantilla-almac} se muestra la plantilla de análisis DAR
aplicada a las herramientas de almacenamiento.

\begin{table}[H]
\centering\footnotesize
\begin{tabular}{|>{\raggedright\arraybackslash\bfseries}p{2.7cm}|c|c|c|c|c|c|c|c|c|}
\hline
\rowcolor{gridhead} \thead[l]{Herramienta} & \thead{Frec.\\SMS} & \thead{Rend.\\(25\%)} & \thead{Esc.\\(20\%)} & \thead{Sop.\\(20\%)} & \thead{Adop.\\(10\%)} & \thead{Obs.\\(10\%)} & \thead{Adec.\\(10\%)} & \thead{Riesgo\\(5\%)} & \thead{DAR} \\ \hline
Ceph & 0 &  &  &  &  &  &  &  &  \\ \hline
\rowcolor{gridrow} Lustre & 0 &  &  &  &  &  &  &  &  \\ \hline
JuiceFS & 0 &  &  &  &  &  &  &  &  \\ \hline
\rowcolor{gridrow} DAOS & 0 &  &  &  &  &  &  &  &  \\ \hline
SeaweedFS & 0 &  &  &  &  &  &  &  &  \\ \hline
\rowcolor{gridrow} ROOT / RDataFrame / Spark / Dask / OSCAR & 1 &  &  &  &  &  &  &  &  \\ \hline
HDFS (EC + ISA-L / SYCL) & 1 &  &  &  &  &  &  &  &  \\ \hline
\rowcolor{gridrow} Knative + x264 / AV1 & 1 &  &  &  &  &  &  &  &  \\ \hline
Longhorn & 0 &  &  &  &  &  &  &  &  \\ \hline
\rowcolor{gridrow} ExaHDF5 & 1 &  &  &  &  &  &  &  &  \\ \hline
HDF5 & 1-2 &  &  &  &  &  &  &  &  \\ \hline
\rowcolor{gridrow} MinIO & 0 &  &  &  &  &  &  &  &  \\ \hline
HDFS & 1 &  &  &  &  &  &  &  &  \\ \hline
\rowcolor{gridrow} RustFS & 0 &  &  &  &  &  &  &  &  \\ \hline
BeeOND & 1 &  &  &  &  &  &  &  &  \\ \hline
\rowcolor{gridrow} OpenStack Swift & 0 &  &  &  &  &  &  &  &  \\ \hline
Garage & 0 &  &  &  &  &  &  &  &  \\ \hline
\rowcolor{gridrow} IPFS & 1 &  &  &  &  &  &  &  &  \\ \hline
GekkoFS & 1 &  &  &  &  &  &  &  &  \\ \hline
\rowcolor{gridrow} GlusterFS & 0 &  &  &  &  &  &  &  &  \\ \hline
BurstFS & 1 &  &  &  &  &  &  &  &  \\ \hline
\rowcolor{gridrow} TEES & 1 &  &  &  &  &  &  &  &  \\ \hline
\end{tabular}
\caption{Plantilla análisis DAR almacenamiento}
\label{tab:plantilla-almac}
\end{table}
